\documentclass[11pt]{article}
\usepackage{icmlko}
\begin{document}

\title{무엇을 세고 있었나 --- 문턱은 내가 골랐고, 실무의 판은 우리가 $59\%$ 로 맞히는 자리에 있으며, 정답은 돈이 아니라 사람 수다}
\author{월드모델 실험실}
\date{노트 528 · 2026년 8월 4일}
\maketitle

\begin{abstract}
노트 527 이 ``접전 $<20$ 점 $\cdot$ 명백 $\ge 50$ 점''으로 갈랐는데
\textbf{그 두 숫자는 내가 골랐다.} 문턱을 방어하는 대신 곡선 전체를 쟀고,
이어서 두 가지를 더 물었다 --- \emph{실무의 판은 그 곡선의 어디에 있나},
그리고 \emph{애초에 무엇을 맞히고 있나}.

\begin{center}
\begin{tabular}{lrr}
\hline
백분위 차 & 쌍 정확도 & 쌍 수 \\
\hline
$0 \sim 5$ 점 & $52.0\%$ & $71{,}791$ \\
$5 \sim 10$ & $54.7\%$ & $73{,}604$ \\
$\mathbf{10 \sim 20}$ & $\mathbf{59.1\%}$ & $134{,}071$ \\
$20 \sim 30$ & $64.7\%$ & $118{,}118$ \\
$50 \sim 60$ & $77.0\%$ & $71{,}168$ \\
$80 \sim 100$ & $89.1\%$ & $31{,}698$ \\
\hline
\end{tabular}
\end{center}

\textbf{문턱이 없다} --- 곡선은 매끈하게 단조 증가한다. 그리고
\textbf{후보 다섯 중 1등과 2등의 실제 백분위 차는 평균 $16.7$ 점}이다.
\textbf{실무에서 마주치는 판은 우리가 $59\%$ 로 맞히는 자리}이지 $81.7\%$
자리가 아니다.
\end{abstract}

\begin{figure}[t]
\centering
\includegraphics[width=\textwidth]{figs/f1.pdf}
\caption{\textbf{가} --- 곡선과 실제 결정의 위치. \textbf{나} --- 결정 단위.
\textbf{다} --- 라벨이 무엇인가.}
\end{figure}

\section{한 노트 만에 내 홍보 문구가 틀렸다}

노트 527 은 ``명백한 판 $81.7\%$'' 를 자랑했다. 그런데 그 구간($\ge 50$ 점)의
쌍은 전체의 $22\%$ 이고, \textbf{실무의 결정은 거기 없다.}

\begin{quote}
후보 다섯을 뽑아 1등과 2등의 실제 성적 차를 재면 평균 $\mathbf{16.7}$ 점이다.
그 구간의 정확도는 $\mathbf{59.1\%}$ 다. \textbf{자랑할 숫자를 고를 때 그
숫자가 실제로 쓰이는 자리에 있는지 먼저 봐야 한다.}
\end{quote}

\section{그래서 결정 단위로 다시 쟀다}

후보 $5$ 개를 무작위로 뽑아 모델 1순위를 고르게 하고, 그것의 \emph{실제 등수}를
$2$ 만 회 $\times$ 씨앗 $4$ 로 셌다.

\begin{center}
\begin{tabular}{lrrr}
\hline
 & 모델 & 무작위 & 배수 \\
\hline
실제 1등을 집는다 & $40.1\%$ & $20.0\%$ & $2.00$ \\
실제 $1\!\sim\!2$ 등 안 & $65.9\%$ & $40.0\%$ & $1.65$ \\
실제 $4\!\sim\!5$ 등을 집는다 & $17.1\%$ & $40.0\%$ & $0.43$ \\
\textbf{실제 꼴찌를 집는다} & $\mathbf{6.2\%}$ & $20.0\%$ & $\mathbf{0.31}$ \\
\hline
\end{tabular}
\end{center}

\begin{quote}
\textbf{모델은 ``제일 좋은 것을 집는 일''보다 ``나쁜 것을 거르는 일''을
훨씬 잘한다.} 1등 적중은 $2.00$ 배인데 꼴찌 회피는 $3.22$ 배다.
제품의 자리는 \emph{최종 선택}이 아니라 \textbf{후보 압축과 사고 방지}다.
\end{quote}

이것은 곡선과 정확히 맞물린다. 꼴찌와 1등은 백분위가 멀어서($\ge 50$ 점)
$77\!\sim\!89\%$ 로 갈리고, 1등과 2등은 가까워서($\approx 17$ 점)
$59\%$ 밖에 안 갈린다.

\section{애초에 무엇을 맞히고 있었나}

이 노트를 쓰게 한 두 번째 물음이다. 코드에서 라벨을 전부 확인했다.

\begin{center}
\begin{tabular}{ll}
\hline
도메인 & 라벨 \\
\hline
팝업 $\cdot$ 시장팝업 & 일평균 \textbf{방문자 수} \\
펀딩 & \textbf{후원자 수} \\
웹툰 & \textbf{관심 등록 수} \\
게임 & \textbf{리뷰 작성자 수}($30$ 일 창) \\
모바일 $\cdot$ 비게임 앱 & \textbf{평가 참여자 수} \\
만화(JP $\cdot$ KR $\cdot$ CN) & \textbf{서재 등록 수} \\
애니 $\cdot$ 세계애니 & \textbf{한줄평 수} \\
아이돌 & \textbf{초동}(첫 주 판매량) \\
도서 & 판매지수(판매량의 함수) \\
\hline
\end{tabular}
\end{center}

\textbf{금액 라벨은 하나도 없다.} 전부 ``충분히 끌려서 실제로 행동한 사람''의
계수다. 코드 주석이 이미 그렇게 적고 있었다 --- 웹툰 라벨 옆에 ``관심 등록 수.
팝업 방문자 $\cdot$ 펀딩 후원자와 같은 물리량이다''.

\begin{quote}
\textbf{왜 매출이 아닌가} --- ① 매출은 기획의 매력이 아니라 \emph{값매김}에
좌우된다(같은 수요도 가격을 바꾸면 매출이 바뀐다). ② 매출은 도메인을 못
건너간다(웹툰 정산과 팝업 정산은 다른 세계다). ③ \textbf{``행동한 사람 수''는
건너간다.} 그래서 웹툰과 팝업이 같은 축 위에 놓인다 ---
\textbf{전이가 되는 뿌리가 라벨 설계에 있다.}
\end{quote}

\textbf{그러므로 못 하는 것도 정해진다} --- ``얼마나 많은 사람이 반응하나''는
맞히고 \textbf{``얼마 버나''는 못 맞힌다.} 방문자가 많아도 객단가가 낮으면
매출은 낮다. 매출을 맞히려면 POS $\cdot$ 정산이 붙어야 하고, 그것이 로드맵의
``방문 $\cdot$ 매출 집계 테이블 $0$ 행'' 항목이다. \textbf{같은 이야기다.}

\section*{조항}

\begin{quote}
\textbf{난이도로 가를 때는 문턱을 고르지 말고 곡선을 낸다.} 문턱은 그 다음에
\emph{쓰임}이 정한다 --- 우리 쓰임은 ``후보 $k$ 개 중 고르기''이고, 그때의
격차 분포가 어느 구간을 봐야 하는지 알려 준다.
\end{quote}

\section*{남은 것}

\begin{itemize}
\item \textbf{관문의 자 다섯이 전부 전체 쌍이다.} 접전 구간이 개선돼도 다섯 자는
거의 안 움직인다 --- \emph{접전 층 전용 자}를 넣을지 정해야 한다.
\item \textbf{도메인별 축 강조}(노트 527 사전 등록) 아직 안 쟀다.
\item \textbf{사람 기준선} --- 사람이 필요하다.
\end{itemize}

\end{document}
